\chapterwithopening{System Testing and Evaluation}{chap:testing_evaluation}{
  \chaptercontentsitem{sec:testing_evaluation_introduction}
  \chaptercontentsitem{sec:testing_strategy}
  \chaptercontentsitem{sec:automated_code_quality_checks}
  \chaptercontentsitem{sec:functional_testing}
  \chaptercontentsitem{sec:role_based_testing}
  \chaptercontentsitem{sec:security_access_control_testing}
  \chaptercontentsitem{sec:non_functional_testing}
  \chaptercontentsitem{sec:preview_deployment_review}
  \chaptercontentsitem{sec:testing_summary}
}

\section{Introduction}\label{sec:testing_evaluation_introduction}

This chapter evaluates EduVerse through the testing evidence visible in the
current repository and the documented development workflow used around that
repository. The evaluation therefore emphasizes verified automated tests,
implemented validation logic, protected-route behavior, and structured manual
review practices rather than unsupported claims of broad test coverage.

The available evidence shows a meaningful testing effort centered on domain
logic, high-risk workflows, access control, and pre-merge review. At the same
time, the repository does not show a dedicated end-to-end browser automation
suite such as Playwright or Cypress. For that reason, the chapter distinguishes
clearly between what was directly evidenced by automated tests, what was
supported by implemented validation logic and review workflow, and what remains
an area for future expansion.

\section{Testing Strategy}\label{sec:testing_strategy}

The testing strategy used for EduVerse can be characterized as a layered
combination of automated logic testing, route-level validation, protected
workflow checks, and manual functional review. The most explicit automated
evidence is the repository test suite executed with \texttt{bun test}, while
additional confidence comes from server-side validation rules, role-aware route
protection, and pre-merge review practices.

This strategy fits the nature of the project. EduVerse contains several
features whose correctness depends on business rules rather than on isolated
screen rendering alone, such as role resolution, class access, exam integrity,
result visibility, feature gating, and browser-safe IDE behavior. Testing those
rules directly at logic and route level helps catch errors in the parts of the
system where academic fairness and access correctness matter most.

At the same time, full cross-screen automation was not identified in the
repository. Consequently, end-user scenarios spanning multiple pages or roles
were supported through manual functional validation and preview-based review
rather than through a complete automated end-to-end suite. This is an important
limitation and is stated explicitly throughout the evaluation.

\section{Automated Code Quality Checks}
\label{sec:automated_code_quality_checks}

Automated testing evidence is clearly present in the repository. The visible
test suite is executed through \texttt{bun test} and focuses on domain logic
and high-risk interaction behavior. In addition, project workflow notes indicate
that GitHub Actions was used to run build, formatting, and type-consistency
checks before merge, thereby extending quality control beyond unit-style tests
alone.

\begin{table}[htbp]
  \centering
  \footnotesize
  \renewcommand{\arraystretch}{1.15}
  \setlength{\tabcolsep}{4pt}
  \caption{Automated Testing and Code Quality Evidence}
  \label{tab:automated_testing_quality_evidence}
  \begin{tabular}{|p{0.11\textwidth}|p{0.25\textwidth}|p{0.18\textwidth}|p{0.29\textwidth}|p{0.09\textwidth}|}
    \hline
    \textbf{Evidence ID} & \textbf{Evidence area} & \textbf{Main role context} & \textbf{Observed repository or workflow evidence} & \textbf{Status} \\
    \hline
    AT-01 & Exam service, grading, and integrity logic & Teacher and Student & Dedicated automated tests cover exam service behavior, scoring logic, and integrity-related rules. & Evidenced \\
    \hline
    AT-02 & Exam lock and manager detail state & Teacher and Student & Test files cover exam lock behavior and manager-side result-state handling. & Evidenced \\
    \hline
    AT-03 & Feature registry resolution & All roles & Automated tests verify feature-registry behavior that determines tool availability. & Evidenced \\
    \hline
    AT-04 & IDE runner and preview logic & Teacher and Student & Automated tests verify browser-safe runner and preview behavior in the coding lab. & Evidenced \\
    \hline
    AT-05 & Class selector and access helpers & Organization Administrator, Teacher, and Student & Tests cover helper logic used to resolve class access and selection behavior. & Evidenced \\
    \hline
    AT-06 & Build, format, and type consistency before merge & Whole project & Project workflow notes indicate GitHub Actions checks for build, formatting, and type consistency. & Evidenced \\
    \hline
  \end{tabular}
\end{table}

\FloatBarrier

The automated evidence is strong for core logic and engineering quality, but it
is narrower than a full acceptance-testing suite. It gives confidence that key
algorithms and helper rules behave correctly, yet it does not by itself prove
every cross-page user journey in a live browser session.

\section{Functional Testing}\label{sec:functional_testing}

Functional testing in EduVerse combined repository-tested logic with manual
validation of user-facing workflows. This was especially important for
scenarios such as organization creation, materials handling, assignment
submission, live participation, and results visibility, where correctness
depends on the interaction of multiple screens, roles, and route handlers.

\begingroup
\footnotesize
\renewcommand{\arraystretch}{1.15}
\setlength{\tabcolsep}{4pt}
\begin{longtable}{|p{0.08\textwidth}|p{0.20\textwidth}|p{0.12\textwidth}|p{0.21\textwidth}|p{0.26\textwidth}|p{0.08\textwidth}|}
  \caption{Representative Functional Testing Scenarios for EduVerse}
  \label{tab:functional_testing_scenarios} \\
  \hline
  \textbf{Test ID} & \textbf{Test Scenario} & \textbf{User Role} & \textbf{Expected Result} & \textbf{Actual Result} & \textbf{Status} \\
  \hline
  \endfirsthead
  \hline
  \textbf{Test ID} & \textbf{Test Scenario} & \textbf{User Role} & \textbf{Expected Result} & \textbf{Actual Result} & \textbf{Status} \\
  \hline
  \endhead
  \hline
  \multicolumn{6}{r}{\small\textit{Continued on the next page}} \\
  \endfoot
  \hline
  \endlastfoot
  FT-01 & Authentication and protected-route entry & Guest/User & Unauthenticated users should be redirected away from protected workspaces. & App-shell protection and route-user enforcement are explicitly implemented. & Evidenced \\
  \hline
  FT-02 & Organization creation with preset feature defaults & Organization Administrator & A new organization should be created with the selected preset and initial feature configuration. & Repository evidence shows organization-creation workflows with presets and feature overrides. & Evidenced \\
  \hline
  FT-03 & Material upload, retrieval, and inline access & Teacher/Student & Valid materials should upload successfully and remain retrievable only through authorized access. & Server-side validation, S3 storage, signed URLs, and inline content routes are implemented. & Evidenced \\
  \hline
  FT-04 & Assignment submission and grading workflow & Teacher/Student & Students should submit within allowed modes, and teachers should review and grade the work. & Submission routes, late rules, file handling, grading, and notification side effects are implemented. & Evidenced \\
  \hline
  FT-05 & Exam start, integrity tracking, grading, and release & Teacher/Student & Eligible students should complete the exam flow while teachers manage review, integrity, and result release. & Exam logic, passcodes, autosave-related behavior, integrity events, grading, retake, flag, and void flows are implemented and partly test-covered. & Evidenced \\
  \hline
  FT-06 & Live-session entry and synchronized whiteboard use & Teacher/Student & Authorized users should join live sessions and collaborate in the shared room environment. & LiveKit token routes, session claiming, realtime updates, and whiteboard synchronization are implemented. & Evidenced \\
  \hline
  FT-07 & Results visibility for students & Organization Administrator/Teacher/Student & Students should see released results only when the class visibility rules permit it. & Results visibility routes and privacy-filtered summary responses are implemented. & Evidenced \\
  \hline
  FT-08 & Browser-based IDE usage & Teacher/Student & Users should open, edit, run, and preview code within the integrated workspace. & IDE interfaces and runner/preview tests confirm the browser-safe workflow; real OS execution is intentionally absent. & Evidenced \\
  \hline
\end{longtable}
\endgroup

These scenarios show that the platform's main academic workflows are strongly
supported by implemented behavior and, in several cases, by direct automated
tests. However, the absence of a full browser-automation suite means that the
project still depended on manual scenario review for integrated user journeys.

\section{Role-Based Testing}\label{sec:role_based_testing}

Because EduVerse is fundamentally role-sensitive, testing role behavior is not
an optional extension to functional testing; it is part of the core system
evaluation. A platform may appear correct at page level while still failing if
roles are resolved incorrectly or if privileged actions leak into student-facing
contexts.

\begin{table}[htbp]
  \centering
  \footnotesize
  \renewcommand{\arraystretch}{1.15}
  \setlength{\tabcolsep}{4pt}
  \caption{Representative Role-Based Testing Scenarios}
  \label{tab:role_based_testing_scenarios}
  \begin{tabular}{|p{0.08\textwidth}|p{0.21\textwidth}|p{0.13\textwidth}|p{0.22\textwidth}|p{0.24\textwidth}|p{0.08\textwidth}|}
    \hline
    \textbf{Test ID} & \textbf{Test Scenario} & \textbf{User Role} & \textbf{Expected Result} & \textbf{Actual Result} & \textbf{Status} \\
    \hline
    RT-01 & Switching the active role inside one organization & Multi-role member & The visible workspace and available actions should change with the selected role. & Selected-role records, role menu behavior, and role-aware API payloads are implemented. & Evidenced \\
    \hline
    RT-02 & Organization administration access & Organization Administrator & Administrative settings and governance pages should remain restricted to the administrative role. & Admin dashboards, organization settings, invite flows, and feature management are role-scoped. & Evidenced \\
    \hline
    RT-03 & Teacher management permissions & Teacher & Teachers should manage only the classes and actions allowed by organization and class rules. & Teacher permission settings and class-management checks are implemented through helpers and RPCs. & Evidenced \\
    \hline
    RT-04 & Student protection from management routes & Student & Students should be blocked from privileged management actions. & Manager-only exam, assignment, and announcement actions are protected in route handlers. & Evidenced \\
    \hline
    RT-05 & Feature gating by organization or class & All roles & Disabled tools should be hidden or blocked consistently for the current context. & Feature-registry and feature-setting logic controls availability at organization and class level. & Evidenced \\
    \hline
  \end{tabular}
\end{table}

\FloatBarrier

The role-based evaluation is one of the strongest parts of the current testing
story because the repository exposes both the underlying authorization model and
the interface patterns that depend on it.

\section{Security and Access Control Testing}
\label{sec:security_access_control_testing}

Security and access control in EduVerse were evaluated primarily through
implementation evidence rather than through a separate penetration-testing
framework. Even so, the repository provides concrete and testable security
mechanisms that can be assessed against expected outcomes.

\begin{table}[htbp]
  \centering
  \footnotesize
  \renewcommand{\arraystretch}{1.15}
  \setlength{\tabcolsep}{4pt}
  \caption{Security and Access Control Testing Scenarios}
  \label{tab:security_access_control_testing_scenarios}
  \begin{tabular}{|p{0.08\textwidth}|p{0.21\textwidth}|p{0.13\textwidth}|p{0.22\textwidth}|p{0.24\textwidth}|p{0.08\textwidth}|}
    \hline
    \textbf{Test ID} & \textbf{Test Scenario} & \textbf{User Role} & \textbf{Expected Result} & \textbf{Actual Result} & \textbf{Status} \\
    \hline
    ST-01 & Unauthorized workspace access & Guest/User & Protected pages and routes should reject unauthenticated access. & App-shell redirects and route-user enforcement are explicitly implemented. & Evidenced \\
    \hline
    ST-02 & Invite acceptance validation & Invited user & Invite completion should depend on the authenticated email identity. & Invite acceptance is bound to the authenticated email rather than token possession alone. & Evidenced \\
    \hline
    ST-03 & Public join-link governance & Guest/User & Join links should respect enablement, expiry, usage limits, and approval rules. & Repository logic supports enable/disable state, expiration, max-use control, default role, and approval-required behavior. & Evidenced \\
    \hline
    ST-04 & Protected file access & Teacher/Student & Downloaded files should remain protected and time-limited. & S3 access is issued through short-lived signed URLs and authenticated content routes. & Evidenced \\
    \hline
    ST-05 & Exam integrity event capture & Student/Teacher & Suspicious exam interruptions should be recorded for later review. & Fullscreen exit, route leave, visibility loss, and blur events are explicitly captured and auditable. & Evidenced \\
    \hline
  \end{tabular}
\end{table}

\FloatBarrier

These results show that security in EduVerse is implemented as layered access
control and verifiable workflow protection rather than as a single gateway
check. The model is appropriate for an educational platform where academic
privacy, role distinction, and assessment fairness all matter.

\section{Non-Functional Testing}\label{sec:non_functional_testing}

Non-functional evaluation in the current project is more limited than
functional evaluation. The available evidence supports maintainability,
security-aware design, and reviewability more strongly than it supports formal
load testing or broad compatibility benchmarking. This limitation does not
invalidate the implementation, but it does define the scope of what can be
claimed.

\begin{table}[htbp]
  \centering
  \footnotesize
  \renewcommand{\arraystretch}{1.15}
  \setlength{\tabcolsep}{4pt}
  \caption{Non-Functional Testing and Evaluation Summary}
  \label{tab:non_functional_testing_summary}
  \begin{tabular}{|p{0.09\textwidth}|p{0.24\textwidth}|p{0.23\textwidth}|p{0.29\textwidth}|p{0.08\textwidth}|}
    \hline
    \textbf{Test ID} & \textbf{Test Scenario} & \textbf{Expected Result} & \textbf{Actual Result} & \textbf{Status} \\
    \hline
    NF-01 & Build, formatting, and type consistency & Project changes should remain buildable and internally consistent before merge. & GitHub Actions checks were documented for build, formatting, and type validation. & Evidenced \\
    \hline
    NF-02 & Maintainability of feature growth & The codebase should support continued refinement without collapsing into one monolithic workflow. & Modular route, feature, and shared-library organization is clearly implemented. & Evidenced \\
    \hline
    NF-03 & Security resilience under normal usage & Protected academic data should remain scoped to authorized users. & Route guards, row-level security, signed file access, and role checks provide strong implementation evidence. & Evidenced \\
    \hline
    NF-04 & Performance under heavier academic usage & Core workflows should remain responsive as classes, files, and messages grow. & Architectural decisions support scalability, but no formal load-testing suite was identified in the repository. & Limited evidence \\
    \hline
    NF-05 & Browser and device compatibility & Main learning workflows should remain usable across common client environments. & The web application is responsive and browser-based, but a formal cross-browser test matrix was not identified. & Limited evidence \\
    \hline
    NF-06 & Usability review before merge & New feature increments should be inspectable before integration. & Preview deployment and review workflow provided a practical means of visual and behavioral inspection. & Evidenced \\
    \hline
  \end{tabular}
\end{table}

\FloatBarrier

The non-functional evaluation therefore supports a balanced conclusion:
EduVerse shows good engineering discipline and modularity, but it would benefit
from stronger formal evidence for load, compatibility, and end-to-end workflow
assurance.

\section{Preview Deployment and Review}\label{sec:preview_deployment_review}

Project development notes show that Vercel preview deployments were used as a
practical review stage before merge. In a system such as EduVerse, this is
particularly valuable because many problems only become visible when a role
switch, feature-gated page, or multi-step workflow is inspected in a realistic
running interface.

When combined with GitHub pull-request review and automated GitHub Actions
checks, preview deployment acts as a safety layer between local implementation
and integration into the main branch. It supports manual validation of visual
behavior, confirms whether a feature increment behaves coherently in the
application shell, and helps the team catch regressions before they reach the
shared codebase.

This review model does not replace a formal end-to-end suite, but it does
provide a defensible intermediate level of evaluation. For EduVerse, it helped
bridge the gap between logic tests and full user-facing academic workflows
during iterative delivery.

\section{Testing Summary}\label{sec:testing_summary}

The evaluation of EduVerse shows a credible testing foundation built on
automated domain tests, route-level validation, role-aware protection, and
preview-based manual review. Core areas such as exam logic, feature gating, IDE
behavior, class access helpers, and security-sensitive routes are supported by
direct repository evidence.

The main limitation is equally clear: a full end-to-end browser automation
suite was not identified, and formal load or compatibility testing evidence is
limited. As a result, the strongest current claims concern correctness of core
logic, access control, and feature behavior, while broader acceptance and scale
assurance remain appropriate targets for future work.
